% !PRINTED 246
\weberpage{246}
%<<CONT
\weq{5}{\Phi_{a,c,\partial} = \frac{\Psi\left[j\left(\frac{c+\partial\omega}{a}\right), j(\omega)\right]}{F'\left[j\left(\frac{c+\partial\omega}{a}\right)\right]}.}

Es gehört also $\Phi_{a,c,\partial}$ dem algebraischen Körper an, der
aus den rationalen Funktionen von $j\left(\frac{c+\partial\omega}{a}\right)$
und $j(\omega)$ besteht, und wir können die Sätze von Bd. I, § 151
anwenden. Aus diesen folgt, daß die $\nu$ Größen $\Phi_{a,c,\partial}$
die Wurzeln einer Gleichung $\nu$ten Grades sind, deren Koeffizienten
rational von $j(\omega)$ abhängen. Wenn die $\nu$ Größen
$\Phi_{a,c,\partial}$ voneinander verschieden sind, so ist diese
Gleichung \gesperrt{irreducibel}. Eine solche Gleichung nennen wir eine zum
\gesperrt{Transformationsgrad $n$ gehörige Transformationsgleichung erster
Stufe}\footnote{In der ersten Auflage habe ich diese Gleichungen
\glqq \gesperrt{invariante Transformationsgleichungen}\grqq{} genannt.
Dieser Ausdruck ist von \gesperrt{Klein} beanstandet worden
(Vorlesungen über ausgewählte Kapitel der Zahlentheorie,
autographiertes Heft, Göttingen 1897). Ich schließe mich
\gesperrt{Kleins} Vorschlag an, indem ich diese Gleichungen jetzt
\glqq \gesperrt{Transformationsgleichungen erster Stufe}\grqq{} nenne,
ohne hier näher auf die Begründung dieses Ausdruckes einzugehen.}.
Jede andere Größe des Körpers kann durch ein solches
$\Phi_{a,c,\partial}$ und durch $j(\omega)$ rational ausgedrückt
werden.

Haben die Funktionen $\Phi_{a,c,\partial}$ die Eigenschaft, für jeden
endlichen Wert von $\omega$ mit positivem, imaginärem Teil, also für
jedes endliche $j(\omega)$ endlich zu bleiben, so ist die Funktion
$\Psi[x,j(\omega)]$ in (4) auch in bezug auf $j(\omega)$ eine
\gesperrt{ganze Funktion}. Die Formel (5) könnte daher nur für solche
\gesperrt{besondere Werte} von $\omega$ versagen, für die zwei Wurzeln der
Invariantengleichung einander gleich werden.

Wenn die $\nu$ Größen $\Phi_{a,c,\partial}$ nicht alle voneinander
verschieden sind, so zerfallen sie in Reihen von gleich vielen
untereinander gleichen, und die aus (5) abzuleitende Gleichung
$\nu$ten Grades ist eine Potenz einer irreducibeln Gleichung, die wir
gelegentlich wohl auch als Transformationsgleichung bezeichnen werden
(Bd. I, § 151, 2).

Sind die Größen $\Phi_{a,c,\partial}$ so gewählt, daß sie für kein
endliches $\omega$ mit positiv imaginärem Bestandteil unendlich
werden, so bleiben sie für jedes endliche $j(\omega)$ endlich, woraus
folgt, daß die Koeffizienten in der Funktion des $\nu$ten Grades
\weq{6}{\Pi(x-\Phi_{a,c,\partial})}
\gesperrt{ganze rationale Funktionen von $j(\omega)$} sind, d.\,h. die
$\Phi_{a,c,\partial}$ sind \gesperrt{ganze algebraische Funktionen von
$j(\omega)$}.
